\documentclass[12pt, a4paper]{article}
\usepackage[utf8]{inputenc}
\usepackage[T1]{fontenc}
\usepackage{amsmath, amssymb, amsthm}
\usepackage{geometry}
\usepackage{enumitem}

% Page layout settings
\geometry{
    margin=1in
}

% Theorem environments
\newtheorem{theorem}{Theorem}
\newtheorem{definition}{Definition}

\title{\textbf{Propositional Logic Encoding for Exam Scheduling}}
\author{}
\date{}

\begin{document}

\maketitle

\section{Problem Definition}

We are tasked with scheduling a set of exams $E$ into a set of time slots $S$.

\begin{itemize}
    \item \textbf{Exams ($E$):} $\{M, P, C, CS, H\}$ corresponding to Mathematics, Physics, Chemistry, Computer Science, and History.
    \item \textbf{Slots ($S$):} $\{1, 2, 3\}$.
\end{itemize}

\subsection{Variables}
We define the propositional variables $X_i$ for each exam $X \in E$ and slot $i \in S$:
$$ X_i = \text{Exam } X \text{ is scheduled in slot } i $$
The total number of variables is $|E| \times |S| = 5 \times 3 = 15$.

\section{Constraint Derivation}

We seek a propositional formula $\Psi$ such that a satisfying assignment to $\Psi$ corresponds exactly to a valid schedule satisfying all constraints.

\subsection{Fundamental Scheduling Constraints}
Before encoding the specific rules of the problem, we must ensure the variables represent a well-formed function $f: E \to S$.

\subsubsection*{1. Completeness (At least one slot)}
Every exam must be scheduled in at least one slot.
$$ \phi_{\text{complete}} = \bigwedge_{X \in E} (X_1 \lor X_2 \lor X_3) $$

\subsubsection*{2. Uniqueness (At most one slot)}
No exam can be scheduled in more than one slot.
$$ \phi_{\text{unique}} = \bigwedge_{X \in E} \big( \neg(X_1 \land X_2) \land \neg(X_1 \land X_3) \land \neg(X_2 \land X_3) \big) $$

The base validity formula is:
$$ \Phi_{\text{base}} = \phi_{\text{complete}} \land \phi_{\text{unique}} $$

\subsection{Problem-Specific Constraints}

\subsubsection*{Constraint 1: M and P Disjointness}
\textit{``M and P cannot be in the same slot.''} \\
For every slot $i$, it is forbidden that both $M_i$ and $P_i$ are true.
$$ \Phi_1 = \bigwedge_{i=1}^{3} \neg(M_i \land P_i) $$
Using logical equivalence $\neg(A \land B) \equiv (\neg A \lor \neg B)$, this can be written as:
$$ \Phi_1 = (\neg M_1 \lor \neg P_1) \land (\neg M_2 \lor \neg P_2) \land (\neg M_3 \lor \neg P_3) $$

\subsubsection*{Constraint 2: C and CS Disjointness}
\textit{``C must be in a different slot from CS.''} \\
Similarly, $C$ and $CS$ cannot share a slot.
$$ \Phi_2 = \bigwedge_{i=1}^{3} \neg(C_i \land CS_i) $$

\subsubsection*{Constraint 3: Slot 1 Requirement}
\textit{``At least one of M, P, C must be in S1.''} \\
This is a simple disjunction of assignments to slot 1.
$$ \Phi_3 = M_1 \lor P_1 \lor C_1 $$

\subsubsection*{Constraint 4: Conditional Constraint}
\textit{``If H is in S1, then M cannot be in S3.''} \\
This is a material implication $H_1 \implies \neg M_3$.
$$ \Phi_4 = H_1 \to \neg M_3 $$
Converting to Clause Normal Form (CNF):
$$ \Phi_4 \equiv \neg H_1 \lor \neg M_3 $$

\subsubsection*{Constraint 5: Non-Consecutive Heavy Courses}
\textit{``CS and M cannot be in consecutive slots.''} \\
Two slots $i$ and $j$ are consecutive if $|i - j| = 1$. The consecutive pairs are $(1,2)$ and $(2,3)$. We must forbid the configuration where one exam is in slot $k$ and the other is in slot $k+1$.

The forbidden configurations are:
\begin{itemize}
    \item $CS$ in 1, $M$ in 2: $(CS_1 \land M_2)$
    \item $M$ in 1, $CS$ in 2: $(M_1 \land CS_2)$
    \item $CS$ in 2, $M$ in 3: $(CS_2 \land M_3)$
    \item $M$ in 2, $CS$ in 3: $(M_2 \land CS_3)$
\end{itemize}

The formula negates all these possibilities:
$$ \Phi_5 = \neg(CS_1 \land M_2) \land \neg(M_1 \land CS_2) \land \neg(CS_2 \land M_3) \land \neg(M_2 \land CS_3) $$

\section{Final Formula}

The complete encoding $\Psi$ is the conjunction of all derived formulas:

$$ \Psi = \Phi_{\text{base}} \land \Phi_1 \land \Phi_2 \land \Phi_3 \land \Phi_4 \land \Phi_5 $$

Writing this out fully:

\begin{align*}
    \Psi = & \left[ \bigwedge_{X \in E} \left( (X_1 \lor X_2 \lor X_3) \land \bigwedge_{1 \le j < k \le 3} \neg(X_j \land X_k) \right) \right] \quad \text{(Valid Schedule)} \\
    \land & \left[ \bigwedge_{i=1}^{3} \neg(M_i \land P_i) \right] \quad \text{(M, P Separate)} \\
    \land & \left[ \bigwedge_{i=1}^{3} \neg(C_i \land CS_i) \right] \quad \text{(C, CS Separate)} \\
    \land & \left[ M_1 \lor P_1 \lor C_1 \right] \quad \text{(S1 Requirement)} \\
    \land & \left[ \neg H_1 \lor \neg M_3 \right] \quad \text{(H and M Conditional)} \\
    \land & \left[ \neg(CS_1 \land M_2) \land \neg(M_1 \land CS_2) \right] \quad \text{(No Consecutive 1-2)} \\
    \land & \left[ \neg(CS_2 \land M_3) \land \neg(M_2 \land CS_3) \right] \quad \text{(No Consecutive 2-3)}
\end{align*}

Every satisfying assignment of $\Psi$ maps uniquely to a valid exam schedule.

\end{document}